% The Weierstrass function on the real axis: a double pole at every lattice point, even,
% with one valley per period bottoming out at a real root e_1.  (Drawn from a truncated
% lattice sum, plotted on pole-avoiding sub-intervals so values stay on-scale.)
\begin{tikzpicture}[>=Latex]
  \draw[->, black!55, line width=0.6pt] (-1.3,0)--(7.6,0) node[right, font=\small] {$x$};
  \draw[->, black!55, line width=0.6pt] (0,-0.2)--(0,5.8) node[above, font=\small] {$\wp(x)$};
  % vertical asymptotes at the lattice points
  \foreach \p in {0,3,6}{\draw[formrust!55, dashed, line width=0.8pt] (\p,-0.1)--(\p,5.55);}
  % the level of the real root reached at the half-period
  \draw[formgold!72!black, dashed, line width=0.7pt]
    (-1.3,1)--(7.6,1) node[right, font=\small, formgold!58!black] {$e_1$};
  % the graph, drawn on intervals that keep a safe margin from each pole
  \foreach \lo/\hi in {-0.85/-0.45, 0.45/2.55, 3.45/5.55, 6.45/6.85}{
    \draw[formteal!78!black, line width=1.6pt, smooth, samples=70, domain=\lo:\hi]
      plot (\x, {1/((\x)^2) + 1/((\x-3)^2) + 1/((\x+3)^2) + 1/((\x-6)^2)});}
  \node[below, font=\small] at (0,-0.04) {$0$};
  \node[below, font=\small] at (3,-0.04) {$\omega$};
  \node[below, font=\small] at (6,-0.04) {$2\omega$};
  \node[formrust!72!black, font=\small, anchor=south] at (4.6,5.15) {double pole at each lattice point};
\end{tikzpicture}
